\chapter*{Введение}

Большие языковые модели (Large Language Models, LLM) стали одной из наиболее значимых технологий последнего десятилетия, трансформировав подходы к обработке естественного языка, генерации текста, машинному переводу и широкому спектру прикладных задач. Модели семейств GPT~\cite{radford2018gpt,radford2019gpt2,brown2020gpt3,achiam2023gpt4}, LLaMA~\cite{touvron2023llama}, Gemini~\cite{team2023gemini}, DeepSeek~\cite{deepseek2024v2} и GigaChat~\cite{2506.09440} демонстрируют способности, ранее считавшиеся недоступными для автоматических систем: связное рассуждение, следование сложным инструкциям, генерация программного кода и анализ мультимодальных данных.

Вместе с ростом возможностей и масштабов внедрения LLM в критически важные приложения возрастает актуальность исследования их безопасности. Мультимодальные модели, обрабатывающие одновременно текстовые и визуальные данные, открывают дополнительные поверхности атак, включая состязательные возмущения изображений, кросс-модальные атаки взлома и внедрение вредоносных инструкций через визуальный канал. Систематическое изучение уязвимостей и разработка защитных механизмов являются необходимым условием безопасного развёртывания данных технологий.

Ключевая гипотеза работы состоит в том, что безопасность мультимодальных LLM должна обеспечиваться не только фильтрацией входов и выходов, но и контролем информационного пространства, из которого модель получает внешние знания. Поэтому поисковый стек рассматривается как часть защитного контура: он ограничивает доступ модели к курируемым источникам, снижает галлюцинации и уменьшает риск генерации ответов на основе вредоносного или недостоверного контента.

%% =================================================================
\section*{Цель, задачи и научная новизна работы}
\addcontentsline{toc}{section}{Цель, задачи и научная новизна работы}
%% =================================================================

\textbf{Актуальность темы.} Промышленное развёртывание русскоязычных и мультимодальных LLM требует одновременно качественного информационного поиска (для заземления ответов на достоверных источниках) и надёжных механизмов безопасности (для сдерживания атак при сохранении полезности). Оба направления для русского языка существенно менее проработаны, чем для английского, а их интеграция в единый защитный контур в литературе систематически не рассматривалась.

\textbf{Целью} работы является разработка методов информационного поиска и переранжирования для русскоязычных больших языковых моделей, а также исследование их интеграции в конвейер безопасного развёртывания мультимодальных LLM через механизм курируемого безопасного поиска.

Для достижения поставленной цели в работе решаются следующие \textbf{задачи}:
\begin{enumerate}
    \item систематизировать современные подходы к построению текстовых и мультимодальных эмбеддингов, методы переранжирования, парадигмы генерации с извлечением информации (RAG), а также системы обеспечения безопасности и уязвимости мультимодальных моделей;
    \item разработать фреймворк обучения высокопроизводительных русскоязычных текстовых эмбеддингов на основе декодерной большой языковой модели и оценить его на бенчмарке MTEB(rus);
    \item предложить компактный реранкер для уточнения результатов первичного извлечения в двухэтапном конвейере поиска и экспериментально оценить его на русскоязычных IR-бенчмарках;
    \item построить поисковый стек, объединяющий лексический, плотный, гибридный и двухэтапный поиск, и обосновать преимущества RAG в части снижения галлюцинаций и устранения необходимости переобучения при обновлении данных;
    \item разработать двухэтапный конвейер обеспечения безопасности мультимодальных моделей, включающий потоковый классификатор обменов и безопасную модель, дообученную с подкреплением с наградой за безопасность;
    \item интегрировать разработанный поисковый стек в конвейер безопасности в виде инструмента безопасного поиска по курируемому индексу и предложить расширение защитных механизмов на визуальную модальность.
\end{enumerate}

\textbf{Объектом исследования} являются большие языковые и мультимодальные модели в контексте их безопасного развёртывания. \textbf{Предметом исследования} выступают методы информационного поиска (эмбеддинги, переранжирование, генерация с извлечением информации) и механизмы обеспечения безопасности, а также их интеграция в единый конвейер.

\textbf{Научная новизна} работы состоит в следующем:
\begin{enumerate}
    \item предложен иерархический трёхэтапный конвейер инструктивного дообучения декодерной LLM GigaChat для построения русскоязычных текстовых эмбеддингов, занявший первое место на русскоязычном бенчмарке MTEB(rus, v1.1)~\cite{2510.22369};
    \item разработан кросс-энкодер реранкер на основе компактных декодерных моделей GigaChat (0.6B и 3B), адаптирующий декодерную LLM к задаче переранжирования заменой каузального внимания на двунаправленное, агрегацией mean pooling и выделенной классификационной головой;
    \item предложен двухэтапный конвейер обеспечения безопасности, в котором независимый от защищаемой модели потоковый классификатор обменов на основе MoE-модели сочетается с безопасной моделью, дообученной с подкреплением с гибридной наградой за безопасность;
    \item впервые поисковый стек на основе эмбеддингов и переранжирования интегрирован в конвейер безопасности в виде инструмента безопасного поиска по курируемому индексу, обеспечивающего защищённой модели актуальный и достоверный взгляд на предметную область, а защитные механизмы распространены на визуальную модальность.
\end{enumerate}

\textbf{Практическая значимость} работы определяется применимостью предложенных методов в промышленных русскоязычных системах: модель эмбеддингов GigaEmbeddings и базовые модели семейства GigaChat выпущены в открытый доступ~\cite{2506.09440,2510.22369,gigaembeddings_hf,gigachat3_10b}, а предложенный конвейер обеспечения безопасности рассчитан на развёртывание с низкими накладными расходами и независимое обновление защитного слоя относительно защищаемой модели.

На \textbf{защиту выносятся} следующие положения:
\begin{enumerate}
    \item адаптация декодерной LLM к двунаправленному вниманию с латентным пулингом и трёхэтапным инструктивным дообучением позволяет получать русскоязычные эмбеддинги уровня state-of-the-art при более чем вдвое меньшем числе параметров ($\approx$3 млрд против 7--8 млрд у сильнейших базовых моделей);
    \item адаптация компактной декодерной LLM к задаче переранжирования (замена каузального внимания на двунаправленное, mean pooling, выделенная классификационная голова) обеспечивает качество уровня лучших открытых реранкеров сопоставимого размера на русскоязычных IR-бенчмарках;
    \item классификация на уровне обмена с пороговым правилом решения и независимостью от защищаемой модели обеспечивает эффективную модерацию контента, а гибридная награда за безопасность позволяет сохранять баланс между безопасностью и полезностью при обучении с подкреплением;
    \item интеграция поискового стека в конвейер безопасности в виде инструмента безопасного поиска по курируемому индексу объединяет преимущества RAG (достоверность, актуальность) с механизмами обеспечения безопасности в единой системе.
\end{enumerate}

\textbf{Личный вклад автора.} Семейство моделей GigaChat (Глава~\ref{ch:gigachat}) разработано коллективом авторов~\cite{2506.09440}; в настоящей работе оно используется как методологическая основа --- источник предобученных бэкбонов для последующих компонентов --- и не выносится на защиту в качестве самостоятельного результата. Фреймворк текстовых эмбеддингов GigaEmbeddings (Глава~\ref{ch:gigaembeddings}) разработан автором: архитектура, трёхэтапный конвейер обучения и экспериментальная оценка~\cite{2510.22369}. Кросс-энкодер реранкер (Глава~\ref{ch:reranking}) предложен и исследован в магистерской диссертации Е.\,А.~Краснопёрова, выполненной под со-руководством автора; автору принадлежат постановка задачи, выбор архитектурного направления и научное руководство, экспериментальная реализация выполнена магистрантом. Поисковый стек (Глава~\ref{ch:rag}) и двухэтапный конвейер обеспечения безопасности (Глава~\ref{ch:safety}), включая модель угроз, потоковый классификатор обменов, дизайн награды за безопасность и интеграцию инструмента безопасного поиска, разработаны автором.

\textbf{Апробация результатов.} Основные результаты работы опубликованы в статьях, посвящённых семейству моделей GigaChat~\cite{2506.09440} и фреймворку GigaEmbeddings~\cite{2510.22369}. Модель эмбеддингов и базовые модели семейства GigaChat выпущены в открытый доступ~\cite{gigaembeddings_hf,gigachat3_10b}; модель GigaEmbeddings по состоянию на июнь 2026~г. занимает первое место в публичном рейтинге MTEB(rus, v1.1).

Соответствие решаемых задач главам диссертации и полученным результатам приведено в Таблице~\ref{tab:intro_map}.

\begin{table}[H]
    \centering
    \caption{Соответствие <<задача~--- глава~--- результат>>.}
    \label{tab:intro_map}
    \begin{tabular}{p{0.26\textwidth}p{0.12\textwidth}p{0.50\textwidth}}
        \toprule
        Задача & Глава & Результат \\
        \midrule
        Обзор методов & Глава~\ref{ch:lit_review} & систематизация направлений: эмбеддинги, переранжирование, RAG, безопасность \\
        Эмбеддинги (GigaEmbeddings) & Глава~\ref{ch:gigaembeddings} & первое место на MTEB(rus, v1.1), средний балл 74.16 \\
        Реранкер & Глава~\ref{ch:reranking} & nDCG@10 $0.78$ на RuBQ; лучшее качество среди открытых моделей сопоставимого размера \\
        RAG-стек & Глава~\ref{ch:rag} & лексический + плотный + гибридный (RRF) + двухэтапный поиск \\
        Конвейер безопасности & Глава~\ref{ch:safety} & потоковый классификатор обменов ($\mathrm{AUC}=0.83$ против $0.60$ у классификатора запросов) \\
        Безопасный индекс & Глава~\ref{ch:safety} & инструмент безопасного поиска по курируемому индексу (заземлённые ответы) \\
        \bottomrule
    \end{tabular}
\end{table}

%% =================================================================
\section*{Теоретические основы}
\addcontentsline{toc}{section}{Теоретические основы}
%% =================================================================

Настоящий раздел кратко вводит понятия, на которых строится изложение; детальный обзор методов вынесен в Главу~\ref{ch:lit_review}.

\textbf{Языковое моделирование.} Языковая модель оценивает вероятностное распределение над последовательностями токенов, $P(w_1,\ldots,w_n) = \prod_{i=1}^{n} P(w_i \mid w_1,\ldots,w_{i-1})$. Классические $n$-граммные модели со сглаживанием~\cite{chen1999empirical} ограничены фиксированным контекстом и разреженностью данных. Переход к нейронным моделям~\cite{bengio2003neural} и рекуррентным сетям --- RNN~\cite{mikolov2010recurrent}, LSTM~\cite{hochreiter1997lstm}, GRU~\cite{cho2014gru} --- позволил моделировать длинные зависимости, однако их последовательная природа ограничивала параллелизацию и масштабирование.

\textbf{Архитектура Transformer.} Предложенная Vaswani et al.~\cite{vaswani2017attention} архитектура (Рисунок~\ref{fig:transformer}) отказалась от рекуррентных вычислений в пользу механизма масштабированного скалярного произведения внимания $\mathrm{Attention}(\mathbf{Q}, \mathbf{K}, \mathbf{V}) = \mathrm{softmax}(\mathbf{Q}\mathbf{K}^\top / \sqrt{d_k})\,\mathbf{V}$, обеспечив параллелизацию обучения и моделирование зависимостей произвольной дальности. Современные модели применяют вращательное позиционное кодирование (RoPE)~\cite{su2024rope} и сформировали три парадигмы предобучения: энкодерную (BERT~\cite{devlin2019bert}), декодерную (GPT~\cite{radford2018gpt,radford2019gpt2,brown2020gpt3}) и энкодер-декодерную (T5~\cite{raffel2020t5}); среди фронтирных LLM доминирует декодерная архитектура благодаря масштабируемости и способности к обучению в контексте (in-context learning).

\begin{figure}[htbp]
    \centering
    \resizebox{0.55\textwidth}{!}{
        \usetikzlibrary{positioning, fit, shapes.geometric, arrows.meta, shadows, calc}

\begin{tikzpicture}[
    node distance=0.8cm and 1cm,
    layer/.style={
        draw,
        rectangle,
        rounded corners,
        minimum width=3cm,
        minimum height=1cm,
        align=center,
        fill=white,
        drop shadow
    },
    process/.style={
        draw,
        rectangle,
        rounded corners,
        minimum width=2.5cm,
        minimum height=0.8cm,
        align=center,
        fill=blue!10
    },
    arrow/.style={
        ->,
        >=Stealth,
        thick
    },
    addnorm/.style={
        draw,
        rectangle,
        rounded corners,
        minimum width=3cm,
        minimum height=0.6cm,
        fill=yellow!20,
        align=center,
        font=\small
    }
]

% Encoder Stack
\node[font=\bfseries] (input) {Inputs};
\node[process, above=0.5cm of input] (emb) {Input Embedding};
\node[circle, draw, fill=white, above=0.3cm of emb] (plus1) {+};
\node[right=0.5cm of plus1] (pos) {Positional Encoding};

\node[layer, fill=orange!20, above=0.5cm of plus1] (mha) {Multi-Head\\Attention};
\node[addnorm, above=0.3cm of mha] (norm1) {Add \& Norm};

\node[layer, fill=blue!10, above=0.5cm of norm1] (ffn) {Feed Forward};
\node[addnorm, above=0.3cm of ffn] (norm2) {Add \& Norm};

% Decoder Stack
\node[right=4cm of input, font=\bfseries, align=center] (output) {Outputs\\(shifted right)};
\node[process, above=0.5cm of output] (outemb) {Output Embedding};
\node[circle, draw, fill=white, above=0.3cm of outemb] (plus2) {+};
\node[right=0.5cm of plus2] (pos2) {Positional Encoding};

\node[layer, fill=orange!20, above=0.5cm of plus2] (mmha) {Masked\\Multi-Head Attention};
\node[addnorm, above=0.3cm of mmha] (norm3) {Add \& Norm};

\node[layer, fill=orange!20, above=0.5cm of norm3] (mha2) {Multi-Head\\Attention};
\node[addnorm, above=0.3cm of mha2] (norm4) {Add \& Norm};

\node[layer, fill=blue!10, above=0.5cm of norm4] (ffn2) {Feed Forward};
\node[addnorm, above=0.3cm of ffn2] (norm5) {Add \& Norm};

\node[process, fill=green!10, above=0.5cm of norm5] (linear) {Linear};
\node[process, fill=green!20, above=0.3cm of linear] (softmax) {Softmax};
\node[above=0.3cm of softmax, font=\bfseries] (probs) {Output Probabilities};

% Connections Encoder
\draw[arrow] (input) -- (emb);
\draw[arrow] (emb) -- (plus1);
\draw[arrow] (pos) -- (plus1);
\draw[arrow] (plus1) -- (mha);
\draw[arrow] (mha) -- (norm1);
\draw[arrow] (norm1) -- (ffn);
\draw[arrow] (ffn) -- (norm2);

% Connections Decoder
\draw[arrow] (output) -- (outemb);
\draw[arrow] (outemb) -- (plus2);
\draw[arrow] (pos2) -- (plus2);
\draw[arrow] (plus2) -- (mmha);
\draw[arrow] (mmha) -- (norm3);
\draw[arrow] (norm3) -- (mha2);
\draw[arrow] (mha2) -- (norm4);
\draw[arrow] (norm4) -- (ffn2);
\draw[arrow] (ffn2) -- (norm5);
\draw[arrow] (norm5) -- (linear);
\draw[arrow] (linear) -- (softmax);
\draw[arrow] (softmax) -- (probs);

% Cross Attention
\draw[arrow] (norm2.north) -- ++(0,0.2) -| (mha2.west);
\draw[arrow] (norm2.north) -- ++(0,0.2) -| (mha2.west); % Simplified, normally K, V

% Residuals (simplified visualization)
\draw[arrow, dashed] (plus1.west) -- ++(-0.5,0) |- (norm1.west);
\draw[arrow, dashed] (norm1.west) -- ++(-0.5,0) |- (norm2.west);

\draw[arrow, dashed] (plus2.east) -- ++(0.5,0) |- (norm3.east);
\draw[arrow, dashed] (norm3.east) -- ++(0.5,0) |- (norm4.east);
\draw[arrow, dashed] (norm4.east) -- ++(0.5,0) |- (norm5.east);

% Labels
\node[draw, dashed, fit=(mha) (norm2), label=left:N$\times$] (enc_block) {};
\node[draw, dashed, fit=(mmha) (norm5), label=right:N$\times$] (dec_block) {};

\node[below=0.1cm of enc_block, font=\bfseries] {Encoder};
\node[below=0.1cm of dec_block, font=\bfseries] {Decoder};

\end{tikzpicture}

    }
    \caption{Архитектура Transformer.}
    \label{fig:transformer}
\end{figure}

\textbf{Mixture of Experts.} Концепция смеси экспертов (Mixture of Experts, MoE)~\cite{jacobs1991moe,shazeer2017moe} реализует принцип условных вычислений (Рисунок~\ref{fig:moe_architecture}): для каждого токена обучаемый маршрутизатор активирует лишь $k \ll N$ экспертов из $N$, что позволяет существенно увеличить ёмкость модели без пропорционального роста вычислительных затрат. Слои MoE заменяют стандартные блоки прямого распространения (FFN), оставляя слой самовнимания без изменений. Подход масштабирован в GShard~\cite{lepikhin2021gshard} (600 млрд параметров) и Switch Transformers~\cite{fedus2022switch} (ускорение предобучения до 7 раз) и лежит в основе Mixtral~\cite{jiang2024mixtral}, DeepSeek~\cite{deepseek2024v2,dai2024deepseekmoe} и GigaChat~\cite{2506.09440}, подробно рассматриваемого в Главе~\ref{ch:gigachat}. Основными вызовами остаются балансировка нагрузки между экспертами и повышенные требования к памяти.

\begin{figure}[htbp]
    \centering
    \resizebox{0.8\textwidth}{!}{
        \usetikzlibrary{positioning, shapes.geometric, arrows.meta, shadows, calc}

\begin{tikzpicture}[
    node distance=1cm and 1.2cm,
    block/.style={
        draw,
        rectangle,
        rounded corners,
        minimum width=2cm,
        minimum height=1cm,
        align=center,
        drop shadow
    },
    expert/.style={
        draw,
        rectangle,
        rounded corners,
        minimum width=1.8cm,
        minimum height=1cm,
        align=center,
        drop shadow
    },
    arrow/.style={
        ->,
        >=Stealth,
        thick
    },
    activearrow/.style={
        ->,
        >=Stealth,
        line width=1.2pt,
        blue!70
    },
    inactivearrow/.style={
        ->,
        >=Stealth,
        thick,
        densely dashed,
        draw=gray!50
    }
]

% Input
\node[block, fill=gray!10] (input) {$\mathbf{x}$\\(Input Token)};

% Router
\node[block, fill=orange!20, above=1.2cm of input] (router) {Gating Network\\(Router)};

% Experts row
\node[expert, fill=blue!25, above left=1.5cm and 3.5cm of router] (e1) {Expert 1\\(FFN)};
\node[expert, fill=gray!8, draw=gray!40, right=0.8cm of e1] (e2) {\color{gray!60}Expert 2\\(FFN)};
\node[expert, fill=blue!25, right=0.8cm of e2] (e3) {Expert 3\\(FFN)};
\node[right=0.6cm of e3, font=\Large] (dots) {$\cdots$};
\node[expert, fill=gray!8, draw=gray!40, right=0.6cm of dots] (en) {\color{gray!60}Expert $N$\\(FFN)};

% Weighted sum
\node[block, fill=green!15, above=1.5cm of router |- e1.north, xshift=4.35cm] (sum) {$\displaystyle\sum_{i} G(\mathbf{x})_i \cdot E_i(\mathbf{x})$};

% Output
\node[block, fill=gray!10, above=1cm of sum] (output) {$\mathbf{y}$\\(Output)};

% Input to Router
\draw[arrow] (input) -- (router);

% Router to active experts
\draw[activearrow] (router.north) -- ++(0, 0.4) -| (e1.south)
    node[pos=0.75, left, font=\small, text=blue!70] {$g_1$};
\draw[activearrow] (router.north) -- ++(0, 0.4) -| (e3.south)
    node[pos=0.75, left, font=\small, text=blue!70] {$g_3$};

% Router to inactive experts (dashed, no arrowhead)
\draw[densely dashed, gray!50, thick] (router.north) -- ++(0, 0.4) -| (e2.south);
\draw[densely dashed, gray!50, thick] (router.north) -- ++(0, 0.4) -| (en.south);

% Active experts to weighted sum
\draw[activearrow] (e1.north) -- ++(0, 0.5) -| ([xshift=-0.6cm]sum.south);
\draw[activearrow] (e3.north) -- ++(0, 0.8) -| ([xshift=0.6cm]sum.south);

% Weighted sum to output
\draw[arrow] (sum) -- (output);

% Annotations
\node[right=0.5cm of router, font=\small, text=black!70, align=left] (annot) {Top-$k$ selection\\($k = 2$ in this example)};


\end{tikzpicture}

    }
    \caption{Архитектура разреженного слоя Mixture of Experts. Маршрутизатор (gating network) для каждого входного токена выбирает подмножество экспертов (top-$k$); неактивные эксперты показаны пунктиром.}
    \label{fig:moe_architecture}
\end{figure}

\textbf{Современные и мультимодальные модели.} Развитие LLM в 2017--2026~гг. обусловлено ростом объёма данных, вычислительных мощностей и числа параметров. Классические законы масштабирования Kaplan et al.~\cite{kaplan2020scaling} и Hoffmann et al. (Chinchilla)~\cite{hoffmann2022chinchilla} установили степенные зависимости качества от вычислительного бюджета; при обучении современных моделей, однако, на исходные Chinchilla-законы ориентируются всё реже, уступая место законам масштабирования пост-обучения, инференса и вычислений на этапе вывода (test-time compute), а также законам масштабирования архитектуры MoE. Открытые модели LLaMA~\cite{touvron2023llama,touvron2023llama2,grattafiori2024llama3}, Mistral~7B~\cite{jiang2023mistral} (с группированным запросным вниманием, GQA), Qwen~\cite{bai2023qwen} и DeepSeek~\cite{deepseek2024v2} (с латентным вниманием, MLA) обеспечивают воспроизводимость исследований. Мультимодальные LLM (MLLM, Рисунок~\ref{fig:multimodal_llm}) объединяют визуальный энкодер (например, ViT~\cite{dosovitskiy2020vit}), адаптер и декодерную LLM (GPT-4V~\cite{achiam2023gpt4}, Gemini~\cite{team2023gemini}, LLaVA~\cite{liu2024llava}); именно они являются основным объектом исследования, поскольку визуальный канал создаёт принципиально новые поверхности атак. Особую значимость для работы представляют русскоязычные модели --- GigaChat~\cite{2506.09440} и разработанный автором фреймворк текстовых эмбеддингов GigaEmbeddings~\cite{2510.22369}.

\begin{figure}[htbp]
    \centering
    \resizebox{0.8\textwidth}{!}{
        \usetikzlibrary{positioning, shapes.geometric, arrows.meta, shadows, calc}

\begin{tikzpicture}[
    block/.style={
        draw,
        rectangle,
        rounded corners,
        minimum width=2.8cm,
        minimum height=1cm,
        align=center,
        drop shadow
    },
    arrow/.style={
        ->,
        >=Stealth,
        thick
    }
]

% Row 1 (y=0): Image processing path
\node[block, fill=gray!10]    at (0,0)   (image) {Image};
\node[block, fill=orange!15]  at (4.5,0) (venc)  {Vision Encoder\\(ViT / CLIP)};
\node[block, fill=blue!10]    at (9,0)   (proj)  {Projection Layer\\(Adapter / MLP)};

% Row 2 (y=-2.5): Text path
\node[block, fill=gray!10]    at (0,-2.5) (text)   {Text Prompt};
\node[block, fill=blue!10]    at (4.5,-2.5) (tokemb) {Tokenizer +\\Embedding};

% LLM block centered between rows
\node[block, fill=purple!12, minimum width=3.2cm, minimum height=3.2cm]
    at (17.5,-1.25) (llm) {Large Language\\Model (Decoder)};

% Output
\node[block, fill=green!12] at (21.5,-1.25) (output) {Generated\\Response};

% Arrows — image path
\draw[arrow] (image) -- (venc);
\draw[arrow] (venc)  -- (proj);
\draw[arrow] (proj.east) -- (llm.west |- proj.east)
    node[midway, above, font=\small] {Visual Tokens};

% Arrows — text path
\draw[arrow] (text) -- (tokemb);
\draw[arrow] (tokemb.east) -- (llm.west |- tokemb.east)
    node[midway, above, font=\small] {Text Tokens};

% Arrow — output
\draw[arrow] (llm) -- (output);

\end{tikzpicture}

    }
    \caption{Архитектура мультимодальной большой языковой модели.}
    \label{fig:multimodal_llm}
\end{figure}

\textbf{Структура работы.} Глава~\ref{ch:lit_review} содержит обзор литературы, систематизирующий основные направления исследований: текстовые и мультимодальные эмбеддинги, методы переранжирования, системы генерации с извлечением информации (RAG), системы обеспечения безопасности ИИ и уязвимости мультимодальных моделей. Глава~\ref{ch:gigachat} посвящена семейству русскоязычных языковых моделей GigaChat, основанных на архитектуре Mixture of Experts, и вводит архитектурную и экспериментальную основу для последующих компонентов. Глава~\ref{ch:gigaembeddings} представляет фреймворк GigaEmbeddings для обучения высокопроизводительных текстовых эмбеддингов на основе декодерной LLM GigaChat-3B. Глава~\ref{ch:reranking} посвящена переранжированию в двухэтапном конвейере поиска и представляет авторский кросс-энкодер реранкер на основе компактных декодерных моделей GigaChat-0.6B и GigaChat-3B. Глава~\ref{ch:rag} рассматривает генерацию с извлечением информации, включая лексический (BM25), гибридный и двухэтапный поиск. Глава~\ref{ch:safety} объединяет результаты предшествующих глав в двухэтапном конвейере обеспечения безопасности мультимодальных моделей, включающем классификатор на основе GigaChat-10B MoE, безопасную модель, дообученную с подкреплением с наградой за безопасность, и инструмент безопасного поиска по курируемому индексу.
